% !TEX root = ../main.tex
\section{Thuật toán áp dụng}

\subsection{Thuật toán Depth-First Search (DFS) cho cơ chế Flood Fill}

Khi người chơi click vào một ô trống (ô có số 0, tức không có mìn lân cận), trò chơi sẽ tự động mở tất cả các ô trống liền kề và dừng lại tại các ô có chứa số (nằm cạnh mìn). Quá trình này yêu cầu việc duyệt một lưới 2D, có thể được mô hình hóa toán học như một đồ thị vô hướng.

Để thực hiện điều này, chúng tôi áp dụng thuật toán \textbf{Tìm kiếm theo chiều sâu không đệ quy (Iterative Depth-First Search - DFS)} kết hợp với một cấu trúc dữ liệu Ngăn xếp (Stack - LIFO).

\subsubsection{Tại sao dùng Không đệ quy (Iterative) thay vì Đệ quy (Recursive)?}
Ngôn ngữ Python ép buộc một giới hạn đệ quy khá khắt khe để tránh lỗi tràn bộ nhớ (mặc định khoảng 1000 khung lệnh). Trên các bảng có kích thước lớn, ví dụ như tùy chỉnh $30\times30$ (900 ô), một hàm DFS đệ quy sâu có thể làm ứng dụng bị đứng (\texttt{RecursionError}). Việc sử dụng cấu trúc Mảng (Array) như một Stack tự quản lý sẽ né tránh hoàn toàn vấn đề này vì nó bị giới hạn chỉ bởi RAM máy tính.

\subsubsection{Các bước của thuật toán DFS}
\begin{enumerate}
    \item Khi có sự kiện click trái mở một ô trống $(r, c)$.
    \item DFS khởi chạy bằng cách đẩy $(r, c)$ vào \texttt{stack}, đồng thời khởi tạo một danh sách \texttt{visited} (các ô đã duyệt qua) và danh sách \texttt{revealed\_cells} chứa kết quả.
    \item Vòng lặp thuật toán tiếp tục chạy chừng nào \texttt{stack} chưa rỗng:
    \begin{enumerate}
        \item Lấy phần tử $(r, c)$ ở đỉnh stack ra để xét.
        \item Nếu $(r, c)$ đã có trong \texttt{visited}, bỏ qua nó. Nếu chưa, thêm $(r, c)$ vào tập \texttt{visited}.
        \item Nếu trạng thái của ô không phải là "chưa mở" (\texttt{hidden}), bỏ qua.
        \item Đổi trạng thái ô thành "đã mở" (\texttt{revealed}) và lưu lại dữ liệu ô đó.
        \item Nếu giá trị trị thực sự của ô là $0$ (trống), tiến hành xét 8 ô hàng xóm lân cận. Nếu ô hàng xóm nằm trong bảng, chưa bị duyệt qua, và không phải là mìn, đẩy tọa độ của nó vào \texttt{stack}.
    \end{enumerate}
    \item Kết thúc và trả về toàn bộ danh sách \texttt{revealed\_cells} để View cập nhật giao diện tổng cộng 1 lần.
\end{enumerate}

\begin{lstlisting}[language=Python, caption=Mã giả thuật toán DFS Flood Fill (Iterative)]
def dfs_flood_fill(start_row, start_col):
    stack = [(start_row, start_col)]
    visited = set()
    revealed_cells = []
    
    while stack:
        r, c = stack.pop()
        if (r, c) in visited: 
            continue
            
        visited.add((r, c))
        if state[r][c] != "hidden":
            continue
            
        state[r][c] = "revealed"
        revealed_cells.append((r, c))
        
        # Neu la o trong (0) -> mo rong 8 huong
        if board[r][c] == 0:
            for nr, nc in get_8_neighbors(r, c):
                if in_bounds(nr, nc) and board[nr][nc] != -1:
                    stack.append((nr, nc))
                    
    return revealed_cells
\end{lstlisting}

\subsubsection{Độ phức tạp thuật toán}
\begin{itemize}
    \item \textbf{Thời gian (Time Complexity):} Gọi $V$ là số lượng ô trống liên thông an toàn, độ phức tạp thời gian cực đại là $O(V)$ vì tại mỗi ô, ta chỉ kiểm tra tối đa 8 ô hàng xóm hằng số $O(1)$. Trên mặt lý thuyết trên một bảng $R \times C$, trường hợp xấu nhất là $O(R \times C)$.
    \item \textbf{Không gian (Space Complexity):} Stack và biến \texttt{visited} sẽ phình to ra tương đương với số ô an toàn liên thông, do đó cực đại là $O(V) \approx O(R \times C)$.
\end{itemize}

\subsection{Cài đặt AI Solver}

Để nâng cao chất lượng bài làm và đáp ứng yêu cầu bộ môn một cách sáng tạo, một hệ thống Trí tuệ Nhân tạo (\texttt{MinesweeperAI}) đã được tích hợp. Thay vì dùng một bot dò đoán ngẫu nhiên thông thường, AI này tự giải quyết toàn bộ ván đấu thông qua \textbf{Hệ thống suy luận 3 tầng}.

\subsubsection{Tầng 1: Propositional Logic (Suy luận theo Mệnh đề)}
Bot liên tục lập bản đồ các ràng buộc và áp dụng hai quy luật mệnh đề tuyệt đối chính xác cho bất kỳ ô số $V$ nào đã mở, đang có $F$ lá cờ bao quanh và giáp với $H$ ô chưa mở:

\begin{itemize}
    \item \textbf{Quy tắc 1 (Tất cả an toàn - All Safe):} 
    Nếu $V = F$, nghĩa là đã tìm đủ mìn xung quanh ô đó. Suy ra toàn bộ $H$ ô chưa mở còn lại phải an toàn. Do đó, tác tử thực hiện hành động \textbf{mở (reveal)} tất cả $H$ ô đó.
    \item \textbf{Quy tắc 2 (Tất cả là mìn - All Mines):} 
    Nếu $(V - F) = H$, nghĩa là số lượng mìn còn thiếu vừa đúng bằng đúng số ô chưa mở nháp. Suy ra tất cả các ô trống đó đều chứa mìn. Do đó, tác tử thực hiện \textbf{cắm cờ (flag)} toàn bộ.
\end{itemize}
Độ phức tạp cho mỗi vòng lặp quét Tầng 1 là cực thấp $O(R \times C \times 8)$, xử lý nhanh $70\%$ - $80\%$ các thế cờ.

\subsubsection{Tầng 2: Thỏa mãn Ràng buộc (Constraint Satisfaction) \& Backtracking}
Khi bộ quy tắc cơ bản ở Tầng 1 gặp ngõ cụt, AI chuyển sang chế độ giải hệ phương trình ràng buộc.

Mỗi ô số đã mở tạo ra một phương trình ràng buộc lên tập hợp các ô giáp nó $C = \{c_1, c_2, ...\}$. 
Ví dụ: $c_1 + c_2 + c_3 = 2$ (có $2$ quả mìn nằm trong $3$ ô này).

\begin{enumerate}
    \item \textbf{Gom nhóm với Union-Find:} Thuật toán phân tách các ô "tiền tuyến" (Frontier cells) chưa mở thành các cụm (cluster) hoàn toàn độc lập với nhau bằng cấu trúc Disjoint Set (Union-Find). Việc này giúp chia để trị, giảm không gian tìm kiếm.
    \item \textbf{DFS Backtracking:} Với mỗi cụm có kích thước $K \le 15$, AI chạy thuật toán đệ quy quay lui (Sử dụng Stack):
    \begin{itemize}
        \item Cây tìm kiếm lần lượt gán giả định \texttt{True} (Mìn) hoặc \texttt{False} (An toàn) cho mỗi ô trong cụm.
        \item Sau mỗi lượt gán, nó sẽ đối chiếu với tập hợp Constraints hiện tại. Nếu vi phạm (VD: ép xuất hiện 2 quả mìn cạnh một ô số 1), nhánh đó bị cắt tỉa (prune) ngay lập tức.
        \item Sau khi tìm ra tất cả cấu hình hợp lệ của cụm, AI chốt đáp án: Nếu một ô nhận giá trị \texttt{True} trong $100\%$ cấu hình, nó chắc chắn là mìn. Ngược lại, $0\%$ thì hoàn toàn an toàn.
    \end{itemize}
\end{enumerate}
Độ phức tạp thời gian ở mức tồi tệ nhất cho mỗi cụm là $O(2^K)$, tuy nhiên do $K$ bị khống chế ở mức rất nhỏ ($K \le 15$) và thuật toán Pruning loại bỏ sớm các nhánh sai, quá trình chạy thực tế chỉ tốn vài mini-giây.

\subsubsection{Tầng 3: Dự đoán Xác suất (Educated Guess)}
Trong những thế cờ ép buộc người chơi phải đoán (các cấu hình 50/50 kinh điển) hoặc khi trên toàn bảng không tìm ra manh mối nào, AI sẽ phân tích rủi ro dựa trên xác suất cục bộ.

Thuật toán duyệt qua mọi ô \texttt{hidden} nằm ở tiền tuyến, xét tất cả các ô số bao quanh nó và đánh giá:
\[ P(Mine) = \max \left( \frac{V - F}{H} \right) \]

Tác tử so sánh $P(Mine)$ của mọi ô trên tiền tuyến, và so với cả tỷ lệ mìn của các ô mù phía sau $\left(\frac{M_{rem}}{H_{rem}}\right)$ (trong đó $M_{rem}$ là số mìn còn lại, $H_{rem}$ là số ô ẩn còn lại). Quyết định \textbf{mở (reveal)} sẽ rơi vào ô có xác suất dính mìn nhỏ nhất hiện tại. Mặc dù Tầng 3 mang tính "đoán", thuật toán tối ưu hóa đáng kể tỷ lệ sinh tồn so với việc phải click ngẫu nhiên của con người.
